\section{Environnement}
\subsection{Entreprise et activité}
\textbf{RUHENNA} est une activité artisanale portée par \textbf{Sumeyra Solak}. Elle propose des prestations au henné (notamment pour des événements) et des produits associés (cônes, accessoires, kits). Le henné est une teinture d’origine végétale (\textit{Lawsonia inermis}) qui permet de réaliser des motifs temporaires sur la peau, avec une dimension culturelle importante lors de certaines célébrations.

\subsection{Organisation et contraintes}
La structure ne disposant pas d’équipe informatique, les échanges ont été directs : Sumeyra porte la vision métier (offre, contenus, règles de réservation) et le stage vise à livrer un outil administrable. Le besoin a été identifié de manière concrète à la suite d’échanges initiés dans un contexte réel (Sumeyra a réalisé le henné lors de mon mariage), puis formalisé en fonctionnalités et règles simples.

\section{Problématique : contexte et objectifs}
Avant la plateforme, la présence en ligne et la gestion des demandes pouvaient être dispersées (réseaux sociaux, messages, échanges manuels). L’objectif a donc été de concevoir un site qui valorise l’activité (galerie, blog, catalogue) tout en structurant la prise de contact et la réservation, avec un back-office permettant une mise à jour autonome.

\section{Spécifications du stage (cahier des charges)}

\subsection{Exigences fonctionnelles}
Le cahier des charges fonctionnel a été construit à partir des besoins concrets de l’activité :
\begin{itemize}
    \item \textbf{Catalogue (boutique)} : afficher les produits sous forme de \textit{catalogue}, avec une page détail par produit.
    \item \textbf{Catégories} : proposer des catégories ordonnées afin de guider la navigation.
    \item \textbf{Galerie} : exposer des réalisations, avec possibilité de mise en avant.
    \item \textbf{Bannière (hero)} : gérer un carrousel d’images sur la page d’accueil.
    \item \textbf{Réservation} : permettre la sélection d’un service, d’une date et d’un créneau disponible, puis l’envoi d’une demande.
    \item \textbf{Contact} : formulaire structuré, stockage du message en base, consultation côté admin.
    \item \textbf{Blog} : création d’articles, publication/dépublication, lecture publique.
    \item \textbf{Administration} : CRUD complet des entités (produits, services, images, articles, messages, rendez-vous).
\end{itemize}

\subsection{Exigences non-fonctionnelles}
\begin{itemize}
    \item \textbf{Maintenabilité} : code lisible, typé, organisé, et facilement extensible.
    \item \textbf{Performance et UX} : mise en cache implicite, chargements fluides, formulaires guidés.
    \item \textbf{SEO} : pages indexables, métadonnées, données structurées (LocalBusiness).
    \item \textbf{Sécurité} : protection de l’admin, validation des entrées, en-têtes de sécurité.
    \item \textbf{Gestion des médias} : upload sécurisé, optimisation des images.
\end{itemize}

\section{Méthodologie et outils mobilisés}

\subsection{Approche de réalisation}
L’activité RUHENNA étant de petite taille, il n’existait pas de méthodologie projet formalisée (Scrum complet, etc.). Le travail a néanmoins été structuré par une logique itérative :
\begin{enumerate}
    \item \textbf{Cadrage} : clarification des fonctionnalités prioritaires (vitrine, galerie, réservation, admin).
    \item \textbf{Modélisation} : construction du schéma de données (Prisma/PostgreSQL).
    \item \textbf{Implémentation verticale} : livraison d’une fonctionnalité complète (UI + API + DB), puis stabilisation.
    \item \textbf{Validation} : tests manuels et retours d’usage (cohérence des formulaires, lisibilité admin).
\end{enumerate}

\subsection{Montée en compétences}
Un point important du stage est le niveau initial : il y avait \textbf{peu de maîtrise préalable} de React/Next.js et de Prisma. L’apprentissage s’est appuyé sur :
\begin{itemize}
    \item la documentation officielle;
    \item des tutoriels en ligne;
    \item l’expérimentation directe.
\end{itemize}
Cette approche a été déterminante : l’apprentissage a progressé au rythme des problèmes réels.

\subsection{Outils et briques techniques du projet}
Les principaux outils observés dans le code sont :
\begin{itemize}
    \item \textbf{Next.js / React} pour l’interface et le routage \textit{App Router} ;
    \item \textbf{TypeScript} pour le typage et la fiabilité ;
    \item \textbf{Prisma} comme ORM ; \textbf{PostgreSQL} comme SGBD ;
    \item \textbf{Neon} pour héberger PostgreSQL de façon compatible avec le déploiement sur \textbf{Vercel} ;
    \item \textbf{Tailwind CSS} pour le style et la cohérence graphique ;
    \item \textbf{Cloudinary} pour le stockage d’images ;
    \item \textbf{Zod} pour la validation des entrées côté API ;
    \item \textbf{TipTap} pour l’édition des articles ;
    \item \textbf{@dnd-kit} pour le glisser-déposer.
\end{itemize}

\section{Choix technologiques : justification et alternatives}

\subsection{Pourquoi React et Next.js ?}
Le choix de React s’explique par la nécessité de construire deux expériences riches :
(i) un front public (galerie, boutique, blog, réservation), et (ii) un back-office interactif (formulaires, modales, tri, filtres).
Next.js a été retenu pour :
\begin{itemize}
    \item son \textbf{routage par convention} (\texttt{src/app/...}) ;
    \item ses \textbf{Route Handlers} permettant d’intégrer l’API sans serveur séparé ;
    \item le rendu hybride (Server Components / Client Components) utile pour la performance et le SEO ;
    \item l’intégration naturelle avec Vercel pour le déploiement.
\end{itemize}
Ces choix se fondent notamment sur la documentation officielle Next.js \cite{nextjsDocs}.

\subsection{Pourquoi Prisma et PostgreSQL ?}
PostgreSQL apporte un modèle relationnel robuste pour des entités structurées (produits, rendez-vous, disponibilités, messages, articles). Prisma complète cet avantage en fournissant :
\begin{itemize}
    \item un \textbf{schéma unique} (\texttt{schema.prisma}) ;
    \item des \textbf{types TypeScript} générés ;
    \item des migrations et une couche d’accès cohérente, limitant les erreurs de requêtes.
\end{itemize}
Cette approche correspond aux bonnes pratiques décrites par Prisma \cite{prismaDocs} : décrire le modèle de données dans un schéma central, versionner les évolutions via des migrations, puis utiliser le client généré (typé) pour effectuer des opérations CRUD de manière cohérente avec le schéma.

\subsection{Comparatif : autres technologies possibles (front, back, base de données)}
Afin de justifier le choix final, un comparatif synthétique d’alternatives réalistes est proposé. L’objectif n’est pas de « prouver » qu’une technologie est meilleure qu’une autre, mais d’expliquer un choix orienté \textbf{simplicité} : RUHENNA est un \textbf{petit site} (peu de données, pas de volumétrie massive) et la priorité était de livrer une solution stable et facile à maintenir.

\renewcommand{\arraystretch}{1.2}
\setlength{\tabcolsep}{6pt}
\begin{table}[H]
    \centering
    \small
    \caption{Comparatif d’alternatives (Front / Back / Base de données) — sources : documentations officielles \cite{reactDocs,nextjsDocs,vercelDocs,prismaDocs,neonDocs}.}
    \label{tab:tech_comparison}
  % --- À mettre dans ton préambule si ce n'est pas déjà fait ---
% \usepackage{array}
% \usepackage{tabularx}
% \usepackage{booktabs}
% \usepackage{multirow}

\begin{tabularx}{\textwidth}{@{} >{\raggedright\arraybackslash}p{3.0cm} >{\raggedright\arraybackslash}p{4.0cm} X >{\raggedright\arraybackslash}p{4.0cm} @{}}
    \toprule
    \textbf{Couche} & \textbf{Option} & \textbf{Avantages} & \textbf{Limites} \\
    \midrule
    \multirow{8}{*}{\textbf{Frontend}} 
        & Next.~js (React) & Bon SEO (SSR/SSG), routing par dossiers, un seul projet pour pages + API. & Courbe d’apprentissage (Server/Client), conventions spécifiques. \\
        \cmidrule{2-4}
        & Nuxt (Vue) & SSR/SSG intégrés, approche cohérente, Vue plus simple par certains profils. & Stack différente ; moins alignée avec l’écosystème React. \\
        \cmidrule{2-4}
        & Angular & Cadre complet, architecture imposée. & Plus lourd pour un petit site ; complexité initiale élevée. \\
    \midrule
    \multirow{8}{*}{\textbf{Backend}} 
        & Route Handlers Next.~js & API intégrée, un seul dépôt, mise en production simple. & À structurer proprement si le projet grossit (services, tests). \\
        \cmidrule{2-4}
        & Express / NestJS & Séparation claire, patterns backend avancés, tests faciles. & Infrastructure en plus, coût de maintenance plus élevé. \\
        \cmidrule{2-4}
        & Django / FastAPI & Productivité Python, écosystème riche. & Double stack (Python + JS), intégration front plus lourde. \\
    \midrule
    \multirow{8}{*}{\textbf{Base de données}} 
        & PostgreSQL & Contraintes, transactions, relations, robustesse. & Modèle relationnel à concevoir correctement. \\
        \cmidrule{2-4}
        & MySQL & Mature, performant, populaire. & Moins de fonctionnalités avancées, choix moins naturel selon l'hôte. \\
        \cmidrule{2-4}
        & MongoDB (NoSQL) & Flexibilité totale du schéma. & Moins adapté aux contraintes fortes (RDV, règles métier strictes). \\
    \bottomrule
\end{tabularx}
\end{table}
\renewcommand{\arraystretch}{1.0}

\section{Schémas pédagogiques : TypeScript et Prisma}

\subsection{Comprendre TypeScript : typage statique et compilation}
TypeScript apporte une vérification de types avant l’exécution. Les types sont \textbf{effacés} après compilation : l’exécution se fait en JavaScript.
Le schéma de la Figure~\ref{fig:ts_pipeline} est inspiré du fonctionnement décrit dans la documentation TypeScript \cite{typescriptDocs}.

\begin{figure}[H]
    \centering
    \begin{tikzpicture}[
        node distance=10mm,
        box/.style={draw, rounded corners, align=center, inner sep=7pt, text width=0.26\textwidth},
        arrow/.style={-{Stealth[length=2.2mm]}, thick}
    ]
    \node[box] (ts) {\textbf{Code TypeScript}\\\small (\texttt{.ts}, \texttt{.tsx})};
    \node[box, right=of ts] (tsc) {\textbf{Compilateur} \texttt{tsc}\\\small vérifie les types};
    \node[box, right=of tsc] (js) {\textbf{JavaScript}\\\small (\texttt{.js}) exécuté};
    \draw[arrow] (ts) -- node[above]{\small compilation} (tsc);
    \draw[arrow] (tsc) -- node[above]{\small émission} (js);
    \end{tikzpicture}
    \caption{Schéma simplifié du fonctionnement TypeScript : vérification à la compilation, exécution en JavaScript.}
    \label{fig:ts_pipeline}
\end{figure}

\subsection{Comprendre Prisma : schéma, migrations et client}
Prisma s’articule autour d’un schéma déclaratif (\texttt{schema.prisma}) décrivant les modèles. À partir de ce schéma, on obtient :
\begin{itemize}
    \item des \textbf{migrations} qui matérialisent les tables/contraintes dans PostgreSQL ;
    \item un \textbf{client Prisma} typé, utilisé dans le code Next.js.
\end{itemize}
La Figure~\ref{fig:prisma_cycle} synthétise le cycle recommandé dans la documentation Prisma \cite{prismaDocs}.

\begin{figure}[H]
    \centering
    \begin{tikzpicture}[
        node distance=10mm,
        box/.style={draw, rounded corners, align=center, inner sep=7pt, text width=0.28\textwidth},
        arrow/.style={-{Stealth[length=2.2mm]}, thick}
    ]
    \node[box] (schema) {\textbf{Schéma Prisma}\\\small \texttt{prisma/schema.prisma}};
    \node[box, right=of schema] (migrate) {\textbf{Migrations}\\\small \texttt{prisma migrate}};
    \node[box, right=of migrate] (db) {\textbf{PostgreSQL}\\\small tables, index, contraintes};
    \node[box, below=of migrate] (client) {\textbf{Prisma Client}\\\small \texttt{prisma generate}};
    \node[box, right=of client] (app) {\textbf{Application Next.js}\\\small requêtes typées};
    \draw[arrow] (schema) -- (migrate);
    \draw[arrow] (migrate) -- (db);
    \draw[arrow] (schema) -- (client);
    \draw[arrow] (client) -- (app);
    \draw[arrow] (app) -- node[above]{\small SQL via driver} (db);
    \end{tikzpicture}
    \caption{Cycle Prisma : du schéma aux migrations et au client typé.}
    \label{fig:prisma_cycle}
\end{figure}

\section{Architecture logicielle à étendre}
Le projet a été développé \textbf{from scratch} : il n’existait pas d’architecture applicative préalable au stage. En revanche, le choix de Next.js impose une architecture par conventions (dossier \texttt{app/}, layouts, routes API) qui sert de cadre structurel et garantit une cohérence globale. Le chapitre~\ref{ch:solution} détaillera cette architecture concrète à travers la structure du dépôt et les principales routes/pages.
